% =====================================================================
%  Выводы по лабораторной работе.
% =====================================================================

\labpartbreak{Выводы}

В ходе выполнения лабораторной работы на основе практических результатов
лабораторной работы №3.1 разработана методика оценки безопасности кода
из шести воспроизводимых, автоматизируемых критериев (устойчивость к
невалидному вводу, потокобезопасность, криптографическая стойкость,
пределы размера входа, работа с конфиденциальными данными в памяти,
изоляция сбоев). Методика применена для сравнения собственной реализации
из лабораторной работы №3.1 с предоставленным алгоритмом -- отдельно
написанной программой, воспроизводящей типичные ошибки разработчика, не
знакомого с этими критериями. Сравнение показало уверенное и
содержательное разделение: собственная реализация полностью
соответствует пяти критериям из шести, предоставленный алгоритм -- лишь
одному.

Семантически работа показала, что методика оценки безопасности кода
приобретает практическую ценность только тогда, когда каждый её критерий
сформулирован как конкретный, воспроизводимый эксперимент с измеримым
результатом, а не как общее пожелание вроде «код должен быть
безопасным». Показательно, что один и тот же критерий -- остаточное
присутствие конфиденциальных данных в памяти после удаления и сборки
мусора -- дал ненулевой результат для обоих алгоритмов, но методика,
включающая уточнённый (двухрежимный, учитывающий кодировку UTF-16LE)
поиск и последующую идентификацию источника остатка, позволила
корректно различить неустранимую особенность платформы (свой алгоритм) и
устранимую ошибку кода (предоставленный алгоритм) -- без этого уточнения
оба случая выглядели бы одинаково «плохими», хотя по существу требуют
совершенно разных действий: для своего алгоритма -- ничего (граница
возможностей платформы), для предоставленного -- удаления
логирования секрета из кода.

С прагматической точки зрения работа показала, что сравнение по чёткой
методике превращает общее, трудно оспоримое утверждение «моя реализация
безопаснее» в конкретный, проверяемый третьей стороной результат:
7 из 7 против 3 из 7 корректных ответов на некорректный ввод, 0 против 2
потерянных записей из 200 при параллельной записи, работающее против
тривиально обратимого шифрование -- каждое утверждение подкреплено
воспроизводимым экспериментом, а не мнением автора о собственном коде.
Также показано, что методика, однажды созданная для одной пары
реализаций, переносима: ни один из шести критериев не привязан к
конкретному языку, фреймворку или предметной области, и та же методика
в дальнейшем может быть применена к любому новому алгоритму, работающему
с конфиденциальными данными в оперативной памяти.
